\documentclass[originalpaper]{jsaiart}

\Vol{12}
\No{1}
\jtitle{Resonating Computer: 内的世界を扱う人工知能の研究的枠組み}
\jsubtitle{記憶・身体感覚・関係性・自己物語を対象としたデータ記述・介入設計・評価指標}
\etitle{Resonating Computer: A New Direction for AI that Engages Personal Internal Worlds}
\esubtitle{Data Modeling, Intervention Design, and Evaluation Challenges for AI Addressing Personal Memory, Embodied Experience, Relationships, and Personal Narratives}

\author{%
 \name{清水}{紘輔}{Kosuke Shimizu}
 \affiliation{筑波大学 情報学群 情報メディア創成学類}%
     {College of Media Arts, Science and Technology, University of Tsukuba}%
     {shimizu@ai.iit.tsukuba.ac.jp}
}

\begin{keyword}
Resonating Computer
Human Augmentation
Personal Informatics
Affective Computing
Artificial Intelligence
\end{keyword}

\begin{summary}
本稿の目的は、AIが個人の内なる「懐かしさ」や「重要な記憶」を引き出す「ナッジ」として機能し、その体験を増幅させることで、ユーザーが自らの経験を「再解釈」するきっかけを作るシステムについて議論することにある．既存の生成モデルや推薦システムは，ユーザの行動履歴や短期的嗜好に基づく外部コンテンツの最適化に重点を置いてきた．一方で，ユーザの現在の内部状態を計測し，それと「共鳴」する過去の記憶を，深い情動喚起を伴う「再構築」を促す介入をするための技術基盤は整備されていない．本稿はこの未踏領域を「Resonating Computer」と定義し，内的世界と計算機の関係を，記録・共鳴・再構築の循環として捉え直す立場を示す．
本稿は，1995年以降のライフログ，Personal Informatics，情動適応対話，ケアロボティクス，およびHCI/VRにおける身体・感覚提示研究を対象としたScoping Review（Arksey \& O'Malley, 2005; Levac et al., 2010）を行い，そこから得られた実装知と評価手続きを整理した．その結果として，内的世界を構成する四つの変数（記憶エピソード，身体化された感覚，関係性のトレース，共鳴パターン）と，AIがそれらに介入する三つの局面（記録，共鳴，再構築）を提示し，四変数×三局面のマトリクスとして体系化する．この分析から，特に「共鳴」と「再構築」の局面に技術的課題が集中していること，すなわち計測された内部状態に基づき，どの記憶をどの身体感覚（メディア）で，どの強度とタイミングで再構築するかを制御する計算論的・倫理的枠組みが欠落していることを指摘する．さらに本稿は，(i) 縦断的な内部状態と記憶エピソードを記述する記録スキーマ，(ii) 没入感を生む感覚メディアの再構築設計，(iii) 共鳴ループを制御・記録する介入機構，(iv) 体験の質（QoE）と内部状態の変化を扱う評価指標群を，中核的な研究領域として提示する．
\end{summary}


\begin{document}
\maketitle

\section{序論}
人工知能の進展により，文章・画像・音声といった「外部コンテンツ」は瞬時に生成可能となり，生成モデルや推薦システムは高い実用性と社会的影響を伴って広く普及している．これらの技術は，ユーザの外在的な行動ログや嗜好パターンを分析し，短期的な関心に応じた出力を簡易にかつ高速に提示する点で顕著な成果をあげてきた\cite{igarashi2025}．既存のAI技術が「次の行動」の予測と外部コンテンツの消費を最適化することに注力してきたのに対し，本稿ではAIが個人の「ふとした時に立ち上がる懐かしさ・記憶・経験\cite{JUHL2023101545,Becker2024Routledge}」を取り扱い，それらをユーザーに提示し，反芻させ，再構築する過程にいかに介入しうるか，という点にある．

本稿は，人工知能研究における新たな探究領域として，「個人の内的世界を参照・反芻を支援する計算枠組み」について議論する．
その位置づけを明確にするために，関連分野のScoping Review\cite{arksey2005,levac2010}を実施し，内的世界に関わる知見の整理と課題抽出を行った．
内的世界を扱うためには，どのデータを記録し，どのように提示し，どのタイミングで介入し，どの指標で評価するかを段階的に整理する必要がある．本稿は，これらの過程を整理するために，\textbf{記録（Record）}・\textbf{提示（Re-present）}・\textbf{再構築（Reconstruct）}という三つの局面を分析軸として採用する．記録局面では，記憶エピソードや身体的感覚，関係性のトレースなどの情報を多感覚情報として取得する方法を扱う．提示局面では，それらの記録を，身体感覚や人との関係を含めて再び感じられるように提示する手法を検討する．再構築局面では，メディア表現から生じる「思い出す」体験や感情の変化を，本人がどのように受け取り，再構成するかを支援・評価する．これら三局面は循環的に作用し，記録された情報は提示や再構築を通じて意味が更新される．記録された情報は提示や再構築を通じて意味が更新される．これら三局面は循環的に作用し，「過去の体験（記録）」が「現在の状態（共鳴）」に応じて「現実空間に再構築」され続けるインタラクティブな系全体を，\textbf{Resonating Computer}と呼ぶ．

本稿は，対象とする研究領域を俯瞰し，未踏領域や研究上の課題を特定することを目的としたScoping Reviewを行う \cite{arksey2005,levac2010}．本稿では，まず第\ref{sec:framework}節において，提案する「四変数」と「三局面」の相互関係を詳細に解説する．次に，それらの枠組みに基づき，内的世界を対象とした先行研究を三局面と四変数の観点から体系的に整理する．続く分析では，四変数×三局面のマトリクスから浮かび上がる未踏領域や偏在傾向を明らかにし，第\ref{sec:rescomp}章において，それらを統合的に捉える枠組みとしての「Resonating Computer」の概念を定義する．最後に，第\ref{sec:agenda}章にて，本稿が提案する研究領域として，共鳴パターンの検出，再構築可能なメディア設計，介入と制御のログ設計，倫理・ガバナンスを含む課題群と，それらに対応する評価指標を提示する．

\section{内的世界を構成する四変数と三局面}\label{sec:framework}
序論で示した四変数と三局面は，本稿全体の分析フレームである．本節では，四つの変数それぞれについて，それが人の内的世界でどのような役割を持ち，AIがそれを扱おうとすると何が問題になるのかを簡潔に位置づける．あわせて，記録・提示・再構築という三局面との関わりを明らかにする．

\subsection{記憶エピソード（Autobiographical Memory Episodes）}
記憶エピソードとは，特定の時間・場所・登場人物・情動を含む出来事のまとまりである\cite{bartlett1932,conway2000}．人はこのまとまりを手がかりとして，自分の過去をたどり直し，これまで自分を取り巻いてきた環境や根底的な価値観を認識できる．回想やReminisenceの研究では，こうしたエピソードを再び取り出せること自体が安心感や自己確認につながることが示されており\cite{kontos2011}，単なる「何が起きたか」という事実以上の意味をもつ．AIの観点から見ると，問題はこのエピソードが自然に「データの単位」にはなっていないことである．現在のライフログ技術は，位置情報や写真，音声，やり取りした相手などを時系列で大量に保存することはできるが\cite{gemmell2002}，そのかたまりが本人にとって「ひとつの出来事」なのか，ただの連続した記録なのかは区別できない．結果として，後からユーザー自身が「そのときの空気ごともう一度思い出したい」と思ったときに，必要なまとまりでアクセスできないという問題が残る．本稿では，この「エピソードというまとまり」を明示的に扱えるようにすることを，記録・提示・再構築のすべてに共通する前提とみなす．すなわち，（i）記録の段階では，出来事の境界とそのときの情動を，後から取り出せる単位として残すこと，（ii）共鳴の段階では，現在の「共鳴パターン」（後述）と照合するキーとして扱えること，（iii）再構築の段階では，その単位ごとに身体感覚や関係性を伴って没入可能に再構築できることである．

\subsection{身体化された感覚（Embodied Sensory Routines）}
身体化された感覚（身体知\cite{suwa2017}）は，姿勢・呼吸・触覚・運動リズムなど言語化しにくいが記憶想起に影響を与える要素である\cite{casey2000,chaudhuri2014}．しかし，身体化された感覚の表現には個人差が著しい．たとえば，同じ触覚刺激であっても，それを認知し，また参照元とする記憶は人により異なる．したがって，刺激の強度・提示の順序・マルチモーダル提示の組み合わせが想起再現率に与える影響を分析する必要がある．AIが身体化された感覚を扱うためには，観測可能な生理指標やモーションデータに加えて，当事者自身による感覚記述を併行して収集し，両者を結びつけたパーソナライズされた表現学習を設計することが求められると考えられる．

\subsection{関係性のトレース（Relational Traces) }
本稿では，人との関わりの累積を，出来事単位ではなく「自分は誰とどんな関係性を持っていたか」という役割感・つながりの感覚として扱い（間柄的存在\cite{Watsuji1979}），これを関係性のトレースと呼ぶ．この変数は，単なるログ化が難しいという点で特徴的である．関係には他者が含まれるため，誰とのどの記録を残し，再び提示してよいかは自分だけで決められない．また，親しい相手に関する写真や手紙の提示が，状況次第で負担や拒否反応を生むことが報告されており\cite{gaver2004}，テレプレゼンスでも過度な「近さ」の再現がかえって緊張を招く例がある\cite{ogawa2012}。つまり，関係性のトレースは，保存すればよい／再提示すればよいという種類のデータではない。三局面との関係は次の通りである。記録（Record）は，やり取りの痕跡を同意とともに残すこと，共鳴（Resonate）は，現在の「共鳴パターン」が特定の他者との関係性をトリガーとすること，再構築（Reconstruct）は，そのつながりをどの強さ・どの距離感でいま再び立ち上げるかを調整することである。


\subsection{共鳴パターン（Resonance Pattern）}
Resonating Computerの核心は，過去の記憶を現在のユーザの状態と「共鳴（関連付け）」させる点にある．したがって，記録・再構築の対象となる過去の体験（変数1, 2, 3）に対し，それらを現在のユーザーの状況と照合し，適切なタイミングで介入を行うことが重要だ．この変数は，ユーザの内部状態（Internal State）と外部文脈（External Context）から構成される．内部状態とは，生理信号（心拍，呼吸，発汗など），表情，声色，姿勢といった，ユーザの現在の情動や認知負荷を反映する計測可能なデータ群を指す．外部文脈とは，現在地，時刻，周囲の環境音，会話の有無といった状況データである．
AIの観点から見ると，問題はこれらの多種多様な時系列データから，いかにして「記憶の引き金（memory cue）」となる特徴を抽出するかである．従来のContext-Aware System\cite{Gerhard2012}・情動計算論\cite{picard1997}が 特定の状況における次の行動の予測に文脈を利用したのに対し，本稿の枠組みでは，AIは「現在の状況に類似しており，かつユーザーが思い出すと喜びを感じる『過去』の体験」を探索するために現在のユーザーの内部状態とユーザーの置かれた状況理解のために駆動される．
三局面との関係は次の通りである．（i）記録の段階では，この「現在の状態」も，記憶エピソード（変数1）と紐づけて時系列で保存される．（ii）再構築の段階では，この「現在の状態」が入力となり，システムが再構築すべきエピソードを選択する主要なキーとなる．（iii）共鳴の段階では，再構築された体験がユーザの「現在の状態」にフィードバックを与え，状態が変化する（例：驚く，落ち着く）。この変化がさらに次の再構築のトリガーとなり，ループを形成する．


\subsection{三局面との連携}
本稿が扱う三局面（記録・共鳴・再構築）は，四つの変数（記憶エピソード，身体化された感覚，関係性のトレース，共鳴パターン）を循環させるための工学的なループとして機能する。
（i）記録は，過去の「記憶エピソード」「身体感覚」「関係性」を，それが発生した時点の「共鳴パターン」と紐付けて保存する． （ii）共鳴は，現在の「共鳴パターン」（変数4）をAIが計測し，それをキーとして，記録（i）の中から最も強く共鳴する過去の「記憶エピソード」（変数1）を検索・選択するプロセスである\cite{wada2007,kidd2008,sengers2005,mcadams2013}．（iii）再演は，共鳴（ii）によって選択されたエピソードを，それに紐づく「身体感覚」（変数2）や「関係性」（変数3）ごと，現実空間に没入可能に再構築するプロセスである．このループは，「現在の状態（共鳴パターン）が，過去の体験（記録）を，現実への没入（再演）として呼び出す」プロセスを記述するものであり，四変数と三局面は相互に定義し合う． この循環は，能動的推論（active inference）や自由エネルギー原理の説明と整合的に捉えられる\cite{friston2009,hesp2021,parr2022activeinference}．能動的推論の枠組みでは，「思い出す」という事象は，現在の経験と「自分にとって世界や経験がどうあるべきか」という内的モデルとの間に生じるずれ（予測誤差）を，過去の貴重な記憶や経験を参照することによって能動的に解消しようとするプロセスとして記述される\cite{friston2009,parr2022activeinference}。この「現在の経験」と「予測誤差を解消するプロセス」という枠組みは、Neisserらが提唱した生態学的記憶論\cite{Neisser1993Perceived}を援用して解釈できる。Neisserらの議論における、記憶を環境内のアフォーダンス（＝現在の経験の手がかり）との相互作用（＝能動的な参照プロセス）として捉える視点は、この枠組みと整合的である。ロボティクスでは，身体感覚と環境の相互作用から将来の状態を予測し，その予測と合うように姿勢や接触を調整する制御則が提案されている\cite{lanillos2021neuroscience,lanillos2020selfother,pabio2021}。本稿の立場に引き寄せれば，記録は「どのずれが自分にとって意味のある出来事だったか」を残す工程，提示は「そのずれにもう一度安全に向き合えるように状況を再構成する工程」，再構築は「そのずれを自分の物語の中でどう位置づけ直すかを本人が決める工程」と読み替えられる。重要なのは，ずれ（違和感や痛み）をただ最小化することではなく，そのズレの解消プロセス（＝再構築）自体を，制御可能で没入的な体験として設計する点である\cite{hesp2021}。


\section{既存の断片的試みと背景的系譜}\label{sec:related-work}
本節では，内的世界を扱うAIの既存研究を，三局面（記録・提示・再構築）と四変数（記憶エピソード，身体感覚，関係性，自己物語）のどこに位置づくかという観点から再整理する。目的は，個別の技術や実践が，どの局面・どの変数には対応できているのか，逆にどこが欠けているのかを可視化し，内的世界の再参照ループ全体に対する課題領域を明らかにすることである。対象とした文献領域は，(i) ライフログとPersonal Informatics，(ii) 生理信号に基づく情動推定と対話適応（Affective Computing），(iii) ケアロボティクスとリミニッセンス技術，(iv) 推薦システムと大規模モデルにおける個別化AI，(v) 自伝的記憶とナラティブ・アイデンティティ研究である。これらはいずれも，人の内的資源（出来事，感覚，関係，語り）を扱おうとするが，扱い方も到達している局面も大きく異なる。文献はACM Digital Library，IEEE Xplore，PubMed，CiNii Articlesなどから収集し，「life-logging」「personal informatics」「affective computing」「reminiscence」「care robotics」「narrative identity」「personalized AI」などの検索語を用いて，主に2012年から2023年の10年間の研究を中心に抽出した。古典的基盤（例：Memex，Augmentation Research Center，Dynabook）については，歴史的系譜を示すために遡及的に参照する。各研究は，(a) 記録・提示・再構築のいずれかを具体的に扱っていること，(b) 内的世界に関する評価・理論的考察があること，(c) 人と計算的プロセスの協働を前提としていることを基準に選定した．

BushのMemexは，個人の経験を検索・連想可能な外部記憶として保持し，あとから自分でたどり直せる環境を構想した\cite{bush1945}。EngelbartはAugmentation Research Centerにおいて，ハイパーテキスト，共同編集，参照リンクの操作など，人の思考過程を外化・再利用可能にする知性拡張環境を実装し\cite{engelbart1962,engelbart1968}，KayとGoldbergはDynabook構想で，個人が自らの経験を編集し直す「パーソナルメディア」という視点を提示した\cite{kay1977}。これらはいずれも，個人が過去の出来事（記憶エピソード）に再アクセスできるように記録・提示の基盤を整備するという点で，現在の「記録／提示」局面の原型となっている。ただし，経験の意味づけをどう再構築し直すかといった部分は，基本的にユーザ個人の負担に委ねられている。RhodesのRemembrance Agentは，ユーザの現在の文脈に応じて，過去のメモやドキュメントを即時に想起させるシステムとして提案された\cite{rhodes1997}。MyLifeBitsは，音声・写真・文書といったライフログ全般を網羅的に保存・検索することで，「全記録（total recall）」を行うことに挑戦した\cite{gemmell2002}．これらは「記憶エピソード」の時系列的・網羅的な保持と再提示には成功しているが，出来事の境界や意味づけ（これは何だったのか，なぜ自分に重要だったのか）をどのように構造化し，再構築まで進めるかという設計は限定的である。原島らによる知的符号化の研究は，人間の知覚単位に沿って音声や映像を圧縮・記述することで，主観的な身体感覚を機械処理可能な特徴量として抽出することの重要性を提唱した\cite{harashima1996}．

ユーザーの嗜好を捉えるための人工知能の開発においては，協調フィルタリングを起点とする推薦システムが，ユーザの嗜好を行動履歴から推定し，潜在因子モデル，ランキング最適化，さらには深層系列モデルへと発展してきた \cite{adomavicius2005,ricci2011,rendle2009bpr,zhang2019,quadrana2018,he2017ncf}．Context-aware Recommendation Systemでは，時間・場所・社会的属性といった文脈情報を状態変数として取り扱い，知識グラフとの統合によって意味的な補完が図られている \cite{adomavicius2011,wang2022}．一方，対話システムやケアロボティクスの領域では，感情推定やユーザ反応に基づいた動作調整を通じて，応答の個別化が進められている \cite{douglascowie2008,cowie2011,kidd2008,wada2007}．近年では，foundation modelに代表される大規模事前学習モデルの文脈において，ユーザ指向の微調整や長期対話履歴の保持を通じた個別化技術が注目されている \cite{bommasani2021,park2023generative,lee2024personalizing}．しかし，これらの「個別化」は多くの場合，ユーザの短期的な行動予測や反応最適化に焦点化しており，ユーザーの重要視したい体験の再参照・再構築プロセスそのものは対象外である．
ソーシャルロボティクス領域では，身体的な刺激や存在感の再構築が，記憶想起や語りを引き出すことが繰り返し報告されている。たとえばPAROは，触覚的な反応と擬似的な交流を通じて高齢者の回想発話を増やすことが示され\cite{wada2007,shibata2011}，Telenoidは，簡素化された身体形状と音声チャネルによって，遠隔にいる相手との親密なコミュニケーションを支えるだけでなく，特定の「誰か」とのつながりの感覚そのものを呼び起こす可能性が示唆されている\cite{ogawa2012}。また，渡邉らによる被爆者写真の彩色プロジェクトは，AIによって白黒写真に色を与え直すことで個人の記憶と歴史的出来事の語りを再び引き出す実践として報告されている\cite{watanave2013,watanave2019}。これらは「身体化された感覚」と「関係性のトレース」を，提示局面における再構築（re-enactment）として扱うアプローチといえる。石井らによるTeleAbsenceは，この系譜をさらに踏み込み，「すでにそこにいない他者」との関係性を再構成することを目指す\cite{ishii2025teleabsence}。これは，触覚・視線・リズムなど，関係性の痕跡を抽出・再提示し，それをユーザーが「これはあの人だ」と再び位置づけ直すプロセスまで含めて設計しようとする点で，単なる記録や提示を超え，「再構築」局面に直接踏み込んでいる。この意味でTeleAbsenceは，「関係性のトレース」という変数に対して，三局面すべてをまたぐ試みとして現時点で最もラディカルな位置にある。

こうした一連の取り組みは，我々の提示する「四変数×三局面」──特に「（ii）身体化された感覚」および「（iii）関係性のトレース」──の枠組みに対応する実践例と捉えることができる．すなわち，これらの感覚提示メディアは，身体・記憶・語り・感情といった内的資源を「提示」の局面でどのように再構築しうるかを問い直しているといえる．

Affective Computing は，音声・映像・生理信号から情動を推定する技術体系として発展してきた \cite{picard1997,cowie2011}．Sensitive Artificial Listener \cite{douglascowie2008} や RECOLA \cite{ringeval2013} といった大規模コーパス\cite{Lang_2005,Ota_2023}は，連続的な情動推定のための標準的評価環境を提供している．
また，筋電図（EMG）計測とArtificial Neural Network(ANN)を用いて笑顔や笑顔の発現を推定する研究 \cite{suzuki2004emg} など，表情筋活動の微細な変化を説明可能な形で扱う試みも報告されている．さらに，マルチモーダル表現学習を用いた長期的情動変化の追跡や個別最適化を目指した感情推定モデルの開発も進展している．しかし，これらの研究は多くの場合，短期的な刺激反応や分類精度の向上に焦点を当てており，推定された情動変化と，個人固有の記憶エピソードや自己物語の再構成プロセスとの接続は十分に検討されていない．情動推定を「内的世界の再参照」として扱うためには，情動を一過的な状態変数ではなく，記録・提示・再構築を横断する内的ダイナミクスとしてモデリングする必要がある．

Personal Informatics は，データの収集・分析・内省・行動変容に至る一連のプロセスを定義し，記録と振り返りのデザインパターンを整理してきた \cite{li2010,choe2014,epstein2015}．長期的セルフトラッキング研究では，個人単位での縦断的データ収集を通じて生活習慣や健康指標の変化を捉え，ログ取得の継続運用に伴う心理的負担やデータ解釈の支援責務を明らかにしている \cite{karkar2017,quadrana2018}．一方，autobiographical design や autoethnography は，研究者自身や当事者による継続的な記述を通じて高密度の主観データを取得し，時間軸付きの個人モデルを構築する方法論として発展してきた \cite{adams2015,ellis2011,neustaedter2012,suwa2015,kagawa2024}．これらの手法は，主観性や再現性に関する批判を受けつつも \cite{anderson2006,atkinson2006,delamont2009}，出来事・関係・身体反応・自己記述を同一タイムライン上で統合的に記録するという実装知を提示している点で重要である．さらに，リフレクティブデザイン \cite{sengers2005} は，問いの開放性を保ちつつ，記憶回顧の主体性を尊重する介入手法として機能しており，Resonating Computer が目指す 記録・提示・再構築の往還構造を設計するうえで重要な理論的基盤を提供する．

これらの調査から明らかになるのは，記憶エピソードの記録や身体化された感覚の提示に関する実践は豊富である一方で，「再構築」の局面においては，四変数すべてにわたり体系的な研究が不足している点である．たとえば，死別後の記憶を設計的に支援するTeleAbsenceの取り組みは萌芽的に見られるものの，評価手法や運用プロトコルはまだ整備されておらず，再構築局面を横断的に捉えるには至っていない\cite{ishii2025teleabsence}．
したがって，次節ではこの課題領域を三つの技術的課題として定式化し，Resonating Computerが果たすべき設計領域を明確化する．

\section{現状の課題}
四変数×三局面のマッピングに基づく分析からは，既存のAI技術では十分に対応されていない領域が明らかになった．これらは以下の三つの型に分類できる．本節では，それぞれの型が生じる背景と，解決に必要な技術的課題を明確化する．
まず推薦システムの多くは，嗜好の「由来」にあたる記憶エピソードや対人関係といった内的資源を，明示的に取り扱わない．多くのモデルはそれらを暗黙的な潜在変数として圧縮しており，ユーザの生活史や身体反応を踏まえて「意味の源泉」を解釈可能な形で表現する仕組みが欠落している．Foundation Modelsに基づくパーソナライゼーションの試み\cite{bommasani2021,park2023generative,lee2024personalizing}においても，そうした個人の生活文脈を参照した嗜好生成の構造化は限定的である．

したがって，記録局面で取得した記憶エピソードや関係性を，再構築局面での推薦や提示に橋渡しする中間表現レイヤー――すなわち「意味のトレース可能な嗜好表現」が，現状では未整備である．

Affective Computingは，連続値としての情動推定やそのアルゴリズム的処理に関して成熟した知見を持つ\cite{picard1997,cowie2011,douglascowie2008,ringeval2013}．一方，自伝的記憶研究では，出来事が時間・場所・関係性と結びついて構成され，回想時にその内容が再構成される過程が明らかにされている\cite{conway2000,bluck2002,pasupathi2001,rubin2009,mcadams2013,ricoeur1991}．しかし，推定された情動を，どの記憶エピソードや自己物語の転換点を再参照すべきかという判断へ接続するための計算論的枠組みは未整備である．情動は単なるリアクションの指標ではなく，再構築の適切なタイミングや方法を調停するメタ情報として用いる必要がある．特に，情動状態によって記憶の意味解釈や語りの負荷が大きく変容することが報告されている\cite{pasupathi2001,mcadams2013}．不適切なタイミングでの刺激提示は，望ましくない再構築や語りの断絶を引き起こすリスクがある．ゆえに，情動推定は提示・介入の「制御指標」として扱う必要があり，そのリスク判断ロジックの欠如が第2のギャップを形成している．

触覚インタフェースや遠隔存在技術などのロボティクス分野では，身体性や関係性が回想や語りに影響を与えることが示されてきた\cite{wada2007,shibata2011,ogawa2012,kidd2008}．また，Personal Informaticsの分野では，縦断的なセルフトラッキングによって個人の状態を長期的に記録・分析する枠組みが発展している\cite{li2010,choe2014,epstein2015,karkar2017,quadrana2018}．しかし，継続的に収集されたログと，身体的あるいは関係的な刺激の提示を統合し，どのタイミングで何を再構築するかを制御するアルゴリズムは確立されていない．縦断データに基づいて介入戦略を導出し，その効果を再びログに反映させるような「循環的な閉ループ設計」が求められている．言い換えれば，身体性・関係性・内的資源の変化を追跡し，それに応じて提示を調整する運用モデルが未整備であり，これが第三の課題である．


\section{Resonating Computer：内的世界を扱う統合的枠組みの提案}\label{sec:rescomp}
本章では，前章のレビューで抽出された過大を踏まえ，横断的知見を設計要件として整理し，それに応答するResonating Computerの構成要素と設計方針を提示する．前章までのレビューにより，内的世界を対象とするAI研究は，メディア装置・センシング・人工知能による調停といった技術的手段を通じて，記録・提示・再構築という三局面を循環的に構成するという共通構造を有していることが明らかとなった．既存研究群はいずれも，特定のデバイスやアルゴリズムのみに依存せず，データ取得，提示制御，介入プロトコルを連携的に統合している点で共通している．

本稿が提案するResonating Computerは，このズレ調整のプロセスにおいて，ユーザの主導権を計算論的に担保することを第一原理とする．どの記憶をどの強度で，どの感覚チャネルを通じて，どのタイミングで再提示するか，そしてその経験を自己物語のどこに位置づけ直すのかという選択において，その選択の主導権をユーザが能動的に行使するための計算的パートナーである．すなわち自由エネルギー原理が示す「主体は予測誤差を扱い続ける」という認知的必然性を踏まえつつ，その予測誤差の扱い方──たとえば，あえて触れない／他者と共有する／自分の言葉で語り直す──を本人が選べるようにするインタフェースである\cite{hesp2021deeplyfelt,joffily2013emotionalvalence}．

記憶や身体感覚のように主観性の強い内的資源は，外部装置や計算処理を介してはじめて再参照可能な形で記録・提示される．そのため，用者自身の語彙や参照スタイルに基づき，経験を再記述するための計算的道具立て（computational toolkit）をいかに提供するかが設計の要点である．いくつかの研究では，提示後の記憶想起や語りにおける自己参照の頻度や時間参照の変化といった指標が評価に用いられており，これは再記述の主体性を示す間接的な指標として有効である．

本レビューからは，\textbf{(1) メディア装置}，\textbf{(2) センシング}，\textbf{(3) アルゴリズム}の三技術軸が，記録・提示・再構築の各局面を支える要素として抽出された．Resonating Computerは，単一のアルゴリズムやデバイスを指すのではなく，どの変数をどのような形で記録し，どのように提示・対話へ反映させ，どの水準で再構築を可能にするかという設計上の問いを束ねる枠組みである．従来のAIが外部環境への最適化やタスク完了を重視してきたのに対し，本枠組みは個人の生活史に基づく再参照可能性を中心に据え，内的世界と技術の響き合いを設計の基点とする．

\subsection{Resonating Computerの中核機能}
\subsubsection{記録と抽出：多感覚情報の意味単位化}
Resonating Computerにおける記録と抽出機能は，記憶エピソード・身体化された感覚・関係性・共鳴パターンといった内的資源を，事後的または準リアルタイムに捉え，後から再参照可能な構造として保存することを目的とする．具体的には，生理信号，音環境，会話，位置情報，対人インタラクションのログを同時に収集し，意味的関連性を保持したまま整理する必要がある\cite{ringeval2013,karkar2017}．さらに，日誌やメモといった自発的記述や身体感覚に関する自己報告を組み合わせることで，データに個人固有の文脈を付与できる\cite{li2010,epstein2015}．しかし，単なるデータ蓄積では内的世界に接続できない．そこで，知的符号化の視点が必要となる．センサ信号や映像を「ユーザーにとって意味を持つ出来事」単位に圧縮する処理が求められ，変化点検知や予測誤差などを指標に，出来事の境界や関係性の転換を検出することができる\cite{harashima1996}．ただし，これらの自動検出は唯一の原理ではなく，本人の語りや感覚的自己報告，心理物理学的な評価\cite{suzuki2004emg}との照合を通じて妥当性を検証する必要がある．さらに，長期的な記録運用では，データ保持に関する倫理的・手続的ガバナンスが不可欠である．ライフログ研究では，削除や非公開化の権利，代理人とのデータ共有などが主要な論点となっており\cite{sellen2010}，Resonating Computerでは，こうした合意情報を記録スキーマの一部として扱い，文脈情報とともにメタデータ化する必要がある．

\subsubsection{提示と再構築：身体化された感覚と関係性の調整}
提示と再構築機能は，記録された素材を再び「触れ直せる」形に構成し直すことで，想起や語りを促す役割を担う．加工された写真や映像が想起を誘発するという報告は，提示が過度に決定的でないほうが効果的である可能性を示唆している\cite{watanave2019}．HodgesらのSenseCam\cite{hodges2006}は，日常の自動撮影によって得られた映像を再生することで，記憶障害者の想起能力を回復させることを示唆した．この結果は，情報提示の抽象度と能動的な記憶回顧のバランスが提示設計において極めて重要であることを示している．また，Woodsらのレビュー\cite{woods2018}は，写真や音声による回想法が感情表現や社会的交流を改善することを報告しており，提示の方法・頻度・強度を文脈に応じて調整する必要性を裏付けている．さらに，自己表象の変化に着目した研究は，自己顔やアバターが行動・態度を変化させる「Proteus効果」を示している\cite{yee2007}．これを踏まえ，自己像を扱う提示設計では，同一化を強制しない粒度・置換比率・出現タイミングの制御が求められる．嗅覚や音圧，呼吸リズムといった身体的情報の提示も同様に，想起と情動を同期させる感覚設計が必要である\cite{casey2000,chaudhuri2014}．Resonating Computerにおける提示機能の効果評価には，追憶の語りの長さや感情語の多寡といった表層指標だけでなく，個人の記憶への認識の変容や再構築の深度を含む質的評価が必要となる．提示が当人の主体性を保ちながら感情変化を伴うかどうかが，重要な評価軸となる．

\subsubsection{再構築：介入調停と評価指標}
再構築機能は，利用者が自身の経験を思い出し，再構築する過程を安全かつ自発的に支援することを目的とする．Personal InformaticsやReflective Designの研究は，気づきの形成を利用者自身に委ねることが長期的な内省を支えると指摘している\cite{li2010,epstein2015}．AIに求められるのは，記録・提示局面で得た素材を参照しつつ，語りのペースや情動変化を観察し，介入強度を適応的に制御する調停ロジックである．Maesらの協調的知的支援研究\cite{maes1994}が示すように，提示や助言が過剰であれば利用者の主導権は失われる．再構築支援アルゴリズムは，語りの沈黙や感情変化，話題の転換といった非言語的シグナルを入力として，必要なときだけ介入するミニマルなポリシーを実装すべきである．評価指標としては，記憶の語りの自己参照率，時間的整合性，介入後の追記率などが有効と考えられる．Autobiographical designやautoethnographyに基づく長期的な自己観察研究\cite{li2010,epstein2015,karkar2017}は，Resonating Computerが記録・提示・回顧を時間軸上で接続する際の理論的基盤を提供している．また，介入ポリシーの最適化の中にユーザーのフィードバックを効果的に接続することができれば，より直接的なアルゴリズムを組むことのできる可能性がある\cite{amershi2019}．記録機能においては，ライフログや情動推定の分野で豊富なベンチマークが存在する一方で，エピソードの境界や関係性の変化といった意味単位を同定する指標が不足している．提示機能では，ロボティクス\cite{ishiguro2012}やメディアアート\cite{Hirakawa_2015}において身体的感覚や関係性の再構築が試みられてきたが，マルチモーダルな提示強度を統合的に調整する制御理論は未整備である．再構築機能に関しては，Personal Informaticsやリフレクティブデザインが介入の構えを提示してきたが，AIが介入のタイミングやその証跡を記録・評価する共通プロトコルは確立されていない．Resonating Computerは，これら三局面にまたがるデータ設計・提示制御・対話の記録と評価を接続する研究とアジェンダして機能する．

\section{Resonating Computerに基づく研究課題}\label{sec:agenda}
Resonating Computerを実現するには，変数と局面を横断する複数の研究課題が相互に連携しなければならない．本章では，特に優先されるべき四つの課題群を提示する．それぞれは，既存研究の断片的成果を踏まえつつ，本枠組みにおける実装的・倫理的要件に基づいて再構成されている．

まずは，記憶エピソードや関係性の転換といった「意味のある瞬間（meaningful moments）」を，本人の記憶と同期しつつ検出可能な指標として定義・抽出することである．既存研究では，生理信号（例：心拍・皮膚電位），音環境，会話パターン，位置情報などのマルチモーダルな時系列データが用いられているが，これらを内的変化と結びつける汎用的な枠組みは確立されていない\cite{ringeval2013,karkar2017}．情動推定研究\cite{douglascowie2008,cowie2011}では，情動と発話の揺らぎが同期的に現れることが示されており，これを応用することで転換点検出のヒューリスティクスが得られる．この課題では，意味単位を構造化するスキーマの設計とともに，推定された転換点と事後的に報告された出来事の一致率，あるいは再構築セッションで頻出した出来事との照合頻度などを評価指標として設定する必要がある．個人に紐づいた縦断的な多感覚情報を，再提示・再構築に活用可能なかたちで収集・解析・統合する共通手続きを確立する方法も議論の余地がある．文献レビューにより，これまでのプロジェクトはデータ取得プロトコルが分散しており，共有可能な形式や再利用性に乏しいことが明らかとなった．本課題において対象とする入力は，同一個人から中長期（数週間〜数ヶ月）にわたり取得された生活ログ，情動自己報告，関係性イベント記録などの多感覚情報である．出力は，個人に固有の状態変化モデルやエピソードクラスタリング結果であり，それらを再提示や再構築に接続するためのインタフェース仕様も含まれる．このような個人ベースの記録と解釈の循環的モデルを整備するには，「一人称研究（first-person research）\cite{suwa2015}」の方法論を手続きとして具体化し，長期的データの取得・解析・共有・再提示までを一貫して扱える体系を構築する必要がある．

続いて，再提示時におけるメディアのタイミング・強度・表現粒度を，感覚的・関係的な記憶の再構築と結び付けて制御する設計指針を構築することも重要だ．レビューからは，多感覚刺激や映像アートといった個別事例は存在するが，提示設計を汎用的に体系化する理論は存在しないことが明らかになった．入力には，エピソードのメタデータ，関係性のタイプ，身体状態のログ，利用者が設定した優先度などが含まれ，出力は提示シナリオ（使用するメディア，提示順序，刺激強度）およびその実行ログである．提示は，忠実な再現だけでなく，記憶想起を誘発する余白や感情的安全性を確保する「設計としてのあいまいさ」が求められる．評価指標には，追憶の継続時間，自己参照の頻度，介入後の情動安定度などを複合的に用い，単なる正確想起では捉えられない再構築効果の全体像を測定する必要がある．

また，再参照を支援しつつ，ユーザーの主体性を損なわない介入ポリシーを定義・調整することも重要な観点の一つだ．特に介入の「しなさ」，つまり非介入の判断が重要であり，沈黙や間合いも含めた制御が鍵となる．評価指標としては，セッション前後の自己効力感，記憶の自己帰属率，AIの介入採用率，介入後の時間参照の変化などを総合的に扱い，再参照ループに与える影響の良否を測定する．

技術的な要素の他に，Resonating Computerが扱う内的資源に対するアクセス，介入，再提示を本人中心に管理可能とする数理的な仕様を定義することである．ライフログ研究では消去権や再同意手続きが個別に検討されてきたが，AIによる再提示・再構築へのアクセス権制御は未整備である．入力には，エピソード単位の公開可否，関係者の同意情報，介入履歴，禁止タグや条件付き提示ルールなどが含まれる．出力は，提示可否のポリシーと，それに基づくリスク制御アクション（停止・延期・刺激制限）である．この枠組みは，単なるプライバシー保護を超えて，関係性と物語を他者が編集しすぎない権利を守るという意味で，根本的に倫理的で政治的な設計課題である\cite{ricoeur2004}．評価には，拒否設定が実際に遵守されていた割合，ポリシー更新の即時性，介入回避率，アクセスログの可読性（第三者監査への対応可能性）などを用い，説明可能性と主権の確保を制度レベルで測定することが必要だと考えられる．

従来の行動変容アプリや自己追跡ツールでは，記録やリマインダーの提示だけでは習慣化が続かず，ユーザーの生活文脈と擦り合わせた設計が不可欠であることが指摘されてきた\cite{stawarz2015}．Resonating Computerにおける長期運用においても，記録の継続性や介入タイミングの妥当性を確保するには，モチベーションの変動や認知負荷に対する動的な調整戦略が求められる．また，ライフログ研究において示唆されているように，長期的なデータ記録には技術的基盤のみならず，参加者との信頼関係構築や継続的なフィードバックループの設計も含まれるべきである\cite{quadrana2018}．

そして，Resonating Computerにおけるシステムの成果は，従来の精度や処理効率といった数理的指標では測定しきれない．むしろ，思い出すことによって「どのような内的資源が再参照されたか」「記憶の構造がいかに変容したか」「関係性のトレースが可視化されたか」といった，人の記憶の質的変化を捉える新たな評価指標が必要となる．評価設計においては，リミニッセンス療法やナラティブ心理学における定量・定性指標を参照する\cite{kontos2011,sedikides2015}．特に，自己物語の変化，記憶の構造変容，他者との関係再構築といった点に焦点を当てる必要がある．加えて，専門家・当事者による質的評価（Narrative Validity）をログ化し，再参照によって開かれた記憶の幅を横断的に比較できる評価指標群として整備する．重要なのは，これらの評価指標は単一の目的関数に還元されないという点である．Resonating Computerは，「想起の誘発，記憶回顧の継続，自己参照性の保持，介入リスクの最小化」など多目的最適化問題として設計・評価されなければならない．このため，評価フレームはトレードオフの構造を明示し，設計者・研究者がどの観測指標を重視した制御戦略を採ったのか，その選択過程自体を記録・再分析可能なログ設計も含める必要がある．

\section{結論}
本稿は，人工知能研究の対象として個人の内的世界を再定義し，記憶エピソード，身体化された感覚，関係性のトレース，自己物語の枠組みという四つの変数を明示した．これらは，従来の生成や推薦を中心としたAI研究では十分に扱われてこなかったが，人間が自己を振り返り，意味づけを行う過程において不可欠な構成要素である．提案する「Resonating Computer」は，これらの変数に対し，記録・提示・再構築という三局面を統合的に設計対象とし，外的な最適化ではなく，内的世界との共鳴を促す新たな研究フレームである．その射程は，AIを外部コンテンツの生成装置として捉える従来の枠組みから，既存の経験資源を再構成する計算的パートナーへと再定義する道筋を提示するものである．
今後の展望としては，ライフログ，情動推定，ロボティクス，ナラティブ研究など，既存領域で蓄積されてきた知見を接続・統合する必要がある．とくに，長期縦断的に記述されたライフデータと，それを再構築するインタフェースを結ぶための共通スキーマの策定と，それを用いた小規模の実践が第一歩である．また，本枠組みを社会実装する上では，制御アルゴリズムと同等の重みで，介入の記録可能性と説明可能性を保証するガバナンス設計が求められる．とくに，介入の「タイミング・強度・根拠」を記録し，本人自身が後から点検・再構築できるようなログ設計は，内的世界を扱うAIにおける研究倫理と実装要件の中心に据えるべきである．本稿が提示した「Resonating Computer」の枠組みと今後の研究領域が，個人の経験と語りを尊重する人工知能研究の基盤となることを期待する．

\begin{acknowledgment}
本稿の構想は，日本ソフトウェア科学会インタラクティブシステムとソフトウェア（ISS）研究会およびUser Interface Software and Technologyにおける議論から多くの示唆を得た．また東京大学名誉教授の原島博氏には，議論の焦点を定めるうえで貴重な助言を賜った．あわせて，議論に参加し建設的なコメントを寄せてくださった関係各位に感謝申し上げる．
\end{acknowledgment}

\bibliography{btxsample}
\bibliographystyle{jsai}

\appendix

\section{付録}
本稿で扱ったResonating Computerの概念は，今後の実装検証に向けて継続的に更新する予定である．各局面のプロトタイプ設計と評価手続きは，別稿にて報告する．

\begin{biography}
\profile{m}{清水 紘輔}{清水紘輔，筑波大学情報学群情報メディア創成学類所属．Human-Computer Interaction, Virtual Realityを興味の主軸としつつ，それを取り巻く人間の情動やそれらをセンシング・計算する仕組み，メディア表現に興味を持ち研究活動を行う．}
\end{biography}

\end{document}
